\Unit[Work $W=\int \vec F\cdot d\vec r$; Work--Energy theorem; potential energies (gravity, spring); mechanical energy conservation; power (instantaneous vs average); conservative fields ($F_x=-\,dU/dx$)]{Work, Energy, Power}
\label{unit:energy}

\begin{KnowledgePoint}{Work and the Work--Energy Theorem \Core}
  \KPFormulas
\begin{formulabox}
  Work transfers energy through a force acting over a displacement; the net work equals the change in kinetic energy.
  \[
  \left\{\begin{aligned}
    &\text{For constant }F\text{ and displacement }d\text{ at angle }\theta:\ \ W=Fd\cos\theta,\\
    &\text{Work--Energy: }\Delta K=W_{\text{net}},\\
    &\text{Power: instantaneous }P=\vec F\cdot\vec v,\quad \text{average over }\Delta t:\ \bar P=\dfrac{W}{\Delta t}.
  \end{aligned}\right.
  \]
\end{formulabox}

  \KPHeuristics
\begin{heuristicsbox}
\begin{itemize}[leftmargin=*]
  \item Choose the system first. Account for energy transferred across its boundary and for internal energy changes within it.
  \item Include nonconservative work explicitly (friction, thrust). For rolling without slipping, static friction often does no work on the rolling body.
  \item Check for indications by comparing the endpoints: $K_i+U_i$ vs. $K_f+U_f$.
\end{itemize}
\end{heuristicsbox}

  \KPProblems
\begin{cheatproblem}
  A $10\,\text{kg}$ block is pushed $5\,\text{m}$ horizontally by a constant force $F=30\,\text{N}$ in the direction of motion. Find the work done by that force.
\begin{solutionbox}
  $W=Fd\cos\theta=30\cdot5\cdot\cos0=150\,\text{J}$.
\end{solutionbox}
\end{cheatproblem}
\begin{cheatproblem}
  A net force of $50\,\text{N}$ accelerates a $5\,\text{kg}$ mass from rest over $10\,\text{m}$. Find the final kinetic energy and speed.
\begin{solutionbox}
  $W_{\text{net}}=50\cdot10=500\,\text{J}$. By work-energy theorem, $\Delta K=W_{\text{net}}=500\,\text{J}$. Then $\tfrac12mv^2=500\Rightarrow v=\sqrt{200}=10\sqrt2\,\text{m/s}$.
\end{solutionbox}
\end{cheatproblem}
\end{KnowledgePoint}

\begin{KnowledgePoint}{Potential Energy and Conservation \Core}
  \KPFormulas
\begin{formulabox}
  When nonconservative work is zero, kinetic plus potential energy is conserved.
  \[
  \left\{\begin{aligned}
    &U_g=mgh\ \text{(near-Earth)},\quad U_s=\tfrac12 kx^2,\\
    &\text{Conservation (when nonconservative work zero): }K_i+U_i=K_f+U_f.
  \end{aligned}\right.
  \]
  \end{formulabox}

  \KPHeuristics
  \begin{heuristicsbox}
  Use reference $U=0$ conveniently; only differences matter.
  \end{heuristicsbox}

  \KPProblems
\begin{cheatproblem}
  A $2\,\text{kg}$ block falls $h=5\,\text{m}$ from rest. Find its speed at the bottom using energy conservation.
\begin{solutionbox}
  $K_i+U_i=K_f+U_f\Rightarrow0+mgh=\tfrac12mv^2+0\Rightarrow v=\sqrt{2gh}=\sqrt{2\cdot9.8\cdot5}=\sqrt{98}\approx9.9\,\text{m/s}$.
\end{solutionbox}
\end{cheatproblem}
\begin{cheatproblem}
  A spring ($k=100\,\text{N/m}$) is compressed by $x=0.2\,\text{m}$. Compute the stored elastic potential energy.
\begin{solutionbox}
  $U_s=\tfrac12 kx^2=\tfrac12\cdot100\cdot0.04=2\,\text{J}$.
\end{solutionbox}
\end{cheatproblem}
\end{KnowledgePoint}

\begin{KnowledgePoint}{Conservative Fields \Extension}
  \KPFormulas
\begin{formulabox}
  In one dimension, a conservative force is the negative spatial derivative of its potential energy:
  \[
  F_x=-\,\dfrac{dU}{dx}
  \]
  Path independence in conservative fields; use energy methods when applicable.
  \end{formulabox}

  \KPHeuristics
  \begin{heuristicsbox}
  Use energy methods if forces are conservative. In a simply connected region, a continuously differentiable force field with zero curl is conservative.
  \end{heuristicsbox}

  \KPProblems
\begin{cheatproblem}
  A mass moves in a potential $U(x)=\tfrac12 kx^2$. Find the force acting on it.
\begin{solutionbox}
  $\vec F=-\dfrac{dU}{dx}\,\hat\imath=-kx\,\hat\imath$ (Hooke's law in 1D).
\end{solutionbox}
\end{cheatproblem}
\begin{cheatproblem}
  Show that gravitational potential $U=-\dfrac{GMm}{r}$ gives force $F=-\dfrac{GMm}{r^2}$ (radial).
\begin{solutionbox}
  First compute $\dfrac{dU}{dr}=\dfrac{d}{dr}\Big(-GMm\,r^{-1}\Big)=+\dfrac{GMm}{r^2}$. Therefore the radial force is $\vec F=-\dfrac{dU}{dr}\,\hat{\mathbf r}=-\dfrac{GMm}{r^2}\,\hat{\mathbf r}$ (attractive, inward).
\end{solutionbox}
\end{cheatproblem}
\end{KnowledgePoint}

\UnitProblemLinks{energy}
